\documentclass[11pt,a4paper]{article}
% Wspólny szablon dla kart supermocy (jedna karta = jedna strona A4)
% Użycie w karcie: \documentclass[11pt,a4paper]{article} % Wspólny szablon dla kart supermocy (jedna karta = jedna strona A4)
% Użycie w karcie: \documentclass[11pt,a4paper]{article} % Wspólny szablon dla kart supermocy (jedna karta = jedna strona A4)
% Użycie w karcie: \documentclass[11pt,a4paper]{article} \input{../preambula} \begin{document}

% Silnik tectonic (XeTeX): polskie znaki przez fontspec, nie przez inputenc/T1
\usepackage{fontspec}
\setmainfont{Arial}
\usepackage[polish]{babel}
\usepackage[a4paper,margin=1.5cm]{geometry}
\usepackage{amsmath,amssymb}
\usepackage{tikz}
\usetikzlibrary{arrows.meta,positioning,calc,shapes.geometric}
\usepackage[most]{tcolorbox}
\usepackage{enumitem}
\usepackage{booktabs}
\usepackage{xcolor}

\pagestyle{empty}
\setlength{\parindent}{0pt}
\setlength{\parskip}{0.4em}

\definecolor{cel}{HTML}{1F4E79}
\definecolor{regula}{HTML}{E8F1FA}
\definecolor{pulapka}{HTML}{FDECEA}
\definecolor{okres}{HTML}{C0392B}
\definecolor{skonczony}{HTML}{1E8449}

% Ramka na regułę
\newtcolorbox{regula}[1][]{
  colback=regula, colframe=cel, boxrule=0.8pt, arc=2pt,
  left=6pt, right=6pt, top=4pt, bottom=4pt, title=#1,
  fonttitle=\bfseries, coltitle=white
}

% Ramka na pułapkę z kartkówki
\newtcolorbox{pulapka}[1][]{
  colback=pulapka, colframe=okres, boxrule=0.8pt, arc=2pt,
  left=6pt, right=6pt, top=4pt, bottom=4pt, title=#1,
  fonttitle=\bfseries, coltitle=white
}

% Nagłówek karty: numer, tytuł, cel
\newcommand{\naglowekKarty}[3]{%
  {\LARGE\bfseries\color{cel} Karta #1: #2}\\[2pt]
  {\large Cel: #3}\par
  \vspace{0.3em}\noindent\rule{\linewidth}{1pt}\par
}

% Pole na obliczenia ucznia: ramka z delikatną kratką, argument = wysokość
\newcommand{\miejsce}[1]{%
  \par\vspace{0.2em}\noindent
  \begin{tikzpicture}
    \draw[step=0.5cm, gray!25, very thin] (0,0) grid (\textwidth,#1);
    \draw[gray!60, rounded corners=2pt] (0,0) rectangle (\textwidth,#1);
    \node[anchor=north west, font=\scriptsize, text=gray] at (0.05,#1) {obliczenia};
  \end{tikzpicture}\par}
 \begin{document}

% Silnik tectonic (XeTeX): polskie znaki przez fontspec, nie przez inputenc/T1
\usepackage{fontspec}
\setmainfont{Arial}
\usepackage[polish]{babel}
\usepackage[a4paper,margin=1.5cm]{geometry}
\usepackage{amsmath,amssymb}
\usepackage{tikz}
\usetikzlibrary{arrows.meta,positioning,calc,shapes.geometric}
\usepackage[most]{tcolorbox}
\usepackage{enumitem}
\usepackage{booktabs}
\usepackage{xcolor}

\pagestyle{empty}
\setlength{\parindent}{0pt}
\setlength{\parskip}{0.4em}

\definecolor{cel}{HTML}{1F4E79}
\definecolor{regula}{HTML}{E8F1FA}
\definecolor{pulapka}{HTML}{FDECEA}
\definecolor{okres}{HTML}{C0392B}
\definecolor{skonczony}{HTML}{1E8449}

% Ramka na regułę
\newtcolorbox{regula}[1][]{
  colback=regula, colframe=cel, boxrule=0.8pt, arc=2pt,
  left=6pt, right=6pt, top=4pt, bottom=4pt, title=#1,
  fonttitle=\bfseries, coltitle=white
}

% Ramka na pułapkę z kartkówki
\newtcolorbox{pulapka}[1][]{
  colback=pulapka, colframe=okres, boxrule=0.8pt, arc=2pt,
  left=6pt, right=6pt, top=4pt, bottom=4pt, title=#1,
  fonttitle=\bfseries, coltitle=white
}

% Nagłówek karty: numer, tytuł, cel
\newcommand{\naglowekKarty}[3]{%
  {\LARGE\bfseries\color{cel} Karta #1: #2}\\[2pt]
  {\large Cel: #3}\par
  \vspace{0.3em}\noindent\rule{\linewidth}{1pt}\par
}

% Pole na obliczenia ucznia: ramka z delikatną kratką, argument = wysokość
\newcommand{\miejsce}[1]{%
  \par\vspace{0.2em}\noindent
  \begin{tikzpicture}
    \draw[step=0.5cm, gray!25, very thin] (0,0) grid (\textwidth,#1);
    \draw[gray!60, rounded corners=2pt] (0,0) rectangle (\textwidth,#1);
    \node[anchor=north west, font=\scriptsize, text=gray] at (0.05,#1) {obliczenia};
  \end{tikzpicture}\par}
 \begin{document}

% Silnik tectonic (XeTeX): polskie znaki przez fontspec, nie przez inputenc/T1
\usepackage{fontspec}
\setmainfont{Arial}
\usepackage[polish]{babel}
\usepackage[a4paper,margin=1.5cm]{geometry}
\usepackage{amsmath,amssymb}
\usepackage{tikz}
\usetikzlibrary{arrows.meta,positioning,calc,shapes.geometric}
\usepackage[most]{tcolorbox}
\usepackage{enumitem}
\usepackage{booktabs}
\usepackage{xcolor}

\pagestyle{empty}
\setlength{\parindent}{0pt}
\setlength{\parskip}{0.4em}

\definecolor{cel}{HTML}{1F4E79}
\definecolor{regula}{HTML}{E8F1FA}
\definecolor{pulapka}{HTML}{FDECEA}
\definecolor{okres}{HTML}{C0392B}
\definecolor{skonczony}{HTML}{1E8449}

% Ramka na regułę
\newtcolorbox{regula}[1][]{
  colback=regula, colframe=cel, boxrule=0.8pt, arc=2pt,
  left=6pt, right=6pt, top=4pt, bottom=4pt, title=#1,
  fonttitle=\bfseries, coltitle=white
}

% Ramka na pułapkę z kartkówki
\newtcolorbox{pulapka}[1][]{
  colback=pulapka, colframe=okres, boxrule=0.8pt, arc=2pt,
  left=6pt, right=6pt, top=4pt, bottom=4pt, title=#1,
  fonttitle=\bfseries, coltitle=white
}

% Nagłówek karty: numer, tytuł, cel
\newcommand{\naglowekKarty}[3]{%
  {\LARGE\bfseries\color{cel} Karta #1: #2}\\[2pt]
  {\large Cel: #3}\par
  \vspace{0.3em}\noindent\rule{\linewidth}{1pt}\par
}

% Pole na obliczenia ucznia: ramka z delikatną kratką, argument = wysokość
\newcommand{\miejsce}[1]{%
  \par\vspace{0.2em}\noindent
  \begin{tikzpicture}
    \draw[step=0.5cm, gray!25, very thin] (0,0) grid (\textwidth,#1);
    \draw[gray!60, rounded corners=2pt] (0,0) rectangle (\textwidth,#1);
    \node[anchor=north west, font=\scriptsize, text=gray] at (0.05,#1) {obliczenia};
  \end{tikzpicture}\par}


\begin{document}

\naglowekKarty{4B}{Zaokrąglanie wersja B: pieniądze i połowa}{zaokrąglać, patrząc tylko na to, czy reszta jest mniejsza niż połowa}

\vspace{0.3em}
\begin{regula}[REGUŁA – JEDNO PYTANIE: CZY JEST POWYŻEJ POŁOWY?]
\textbf{Krok 1.} Znajdź dwie liczby zaokrąglone, między którymi leży Twoja liczba, np. $15{,}3$ i $15{,}4$ (krok to $0{,}1$).\\
\textbf{Krok 2.} Policz \textbf{środek}, czyli połowę kroku: $15{,}3+0{,}05=15{,}35$.\\
\textbf{Krok 3.} Porównaj: poniżej środka $\rightarrow$ bliżej mniejszej liczby. Powyżej środka $\rightarrow$ bliżej większej.
\end{regula}

\vspace{0.4em}
\textbf{Druga droga: pieniądze.} Zaokrąglanie do pełnych złotych to pytanie, czy mamy więcej niż $50$ groszy:
$$4{,}62\ \text{zł}=4\ \text{zł}\ 62\ \text{gr}\ \xrightarrow{\ 62>50\ }\ \mathbf{5\ zł}
\qquad\qquad
4{,}38\ \text{zł}=4\ \text{zł}\ 38\ \text{gr}\ \xrightarrow{\ 38<50\ }\ \mathbf{4\ zł}$$

\begin{center}
\begin{tikzpicture}[>=Stealth, x=10cm]
  % Krok 1 zł: od 4 do 5; środek 4,5
  \draw[->] (-0.05,0) -- (1.05,0) node[right, font=\small] {zł};
  \foreach \v/\lab in {0/{4\,zł}, 0.5/{4{,}5\,zł}, 1/{5\,zł}} {
    \draw (\v,0.1) -- (\v,-0.1);
    \node[below, font=\small] at (\v,-0.12) {$\lab$};
  }
  % 4,62 zł: powyżej połowy -> bliżej 5 zł
  \fill[okres] (0.62,0) circle (2.2pt) node[above=4pt, font=\small\bfseries, color=okres] {$4{,}62$};
  \draw[dashed, okres] (0.62,0) -- (0.62,-0.45);
  \node[font=\small, color=okres, anchor=north] at (0.62,-0.5) {$62\,\text{gr}>50\,\text{gr}\Rightarrow 5\,\text{zł}$};
  % 4,38 zł: poniżej połowy -> bliżej 4 zł
  \fill[cel] (0.38,0) circle (2.2pt) node[below=4pt, font=\small\bfseries, color=cel] {$4{,}38$};
  \node[font=\small, color=cel] at (0.38,0.6) {$38\,\text{gr}<50\,\text{gr}\Rightarrow 4\,\text{zł}$};
\end{tikzpicture}
\end{center}

\vspace{0.3em}
\textbf{Przykłady rozstrzygnięte} (zawsze: sąsiedzi, środek, porównanie):
\begin{center}
\renewcommand{\arraystretch}{1.4}
\begin{tabular}{@{}llll@{}}
\toprule
Zadanie & Sąsiedzi & Środek & Wynik \\
\midrule
$4{,}8\,|\,29$ do dziesiątych & $4{,}8$ i $4{,}9$ & $4{,}85$ & $4{,}829<4{,}85\Rightarrow \mathbf{4{,}8}$ \\
$15{,}3\,|\,7$ do dziesiątych & $15{,}3$ i $15{,}4$ & $15{,}35$ & $15{,}37>15{,}35\Rightarrow \mathbf{15{,}4}$ \\
$0{,}351$ do setnych & $0{,}35$ i $0{,}36$ & $0{,}355$ & $0{,}351<0{,}355\Rightarrow \mathbf{0{,}35}$ \\
$7{,}286$ do setnych & $7{,}28$ i $7{,}29$ & $7{,}285$ & $7{,}286>7{,}285\Rightarrow \mathbf{7{,}29}$ \\
$14{,}52$ do jedności & $14$ i $15$ & $14{,}5$ & $14{,}52>14{,}5\Rightarrow \mathbf{15}$ \\
$0{,}(6)=0{,}666\ldots$ do setnych & $0{,}66$ i $0{,}67$ & $0{,}665$ & $0{,}666>0{,}665\Rightarrow \mathbf{0{,}67}$ \\
$3{,}998$ do setnych & $3{,}99$ i $4{,}00$ & $3{,}995$ & $3{,}998>3{,}995\Rightarrow \mathbf{4{,}00}$ \\
\bottomrule
\end{tabular}
\end{center}

\begin{pulapka}[PUŁAPKI Z KARTKÓWKI]
\begin{enumerate}[leftmargin=*, itemsep=2pt]
  \item \textbf{Połowa to połowa kroku, nie całości.} Dla setnych krok to $0{,}01$, więc środek to $0{,}005$ powyżej zaokrąglonej liczby. Dla jedności krok to $1$, więc środek to $0{,}5$.
  \item \textbf{Zera zostają.} Wynik $4{,}00$ ma dwie cyfry po przecinku, bo zaokrąglamy do setnych.
  \item \textbf{Okres najpierw rozpisz:} $0{,}(6)=0{,}666\ldots$. Dopiero wtedy szukasz środka.
\end{enumerate}
\end{pulapka}

\end{document}
